\subsection{End-of-Chapter Problems}

\begin{enumerate}

% Section 1.1 Problems
\item \textbf{Classification of Control Systems}
\begin{enumerate}
\item Classify each of the following systems as open-loop or closed-loop, continuous-time or discrete-time, and linear or nonlinear:
\begin{enumerate}
\item A home thermostat controlling a furnace based on temperature readings
\item A cruise control system in an automobile using only a speedometer
\item An industrial robot arm controlled by pre-programmed joint angles
\item A digital camera's autofocus system that adjusts lens position based on image sharpness
\end{enumerate}
\item For each system in part (a), identify the reference input, control input, measured output, and disturbance signals (if any exist).

\item \textbf{Design Problem:} Sketch the block diagram for a coffee maker that maintains water temperature at a specified setpoint. Identify all sensors, actuators, and controller elements. Is this system open-loop or closed-loop? Explain your reasoning.

\end{enumerate}

% Section 1.2 Problems
\item \textbf{Electrical System Modeling}
\begin{enumerate}
\item Derive the differential equation relating the output voltage $v_o(t)$ to the input voltage $v_i(t)$ for the circuit shown in Figure 1.XX, consisting of a resistor $R$, capacitor $C$, and inductor $L$ in series connected to a load resistor $R_L$.

\item For the circuit in part (a), determine the transfer function $G(s) = V_o(s)/V_i(s)$ assuming zero initial conditions.

\item A passive RC low-pass filter has $R = 10\text{ k}\Omega$ and $C = 1\text{ \mu F}$. Determine the cutoff frequency in radians per second and hertz. If the input is a unit step function, find the expression for the output voltage $v_o(t)$.

\end{enumerate}

\item \textbf{Mechanical System Modeling}
\begin{enumerate}
\item A mass-spring-damper system has mass $m = 2\text{ kg}$, spring constant $k = 100\text{ N/m}$, and damping coefficient $c = 10\text{ N·s/m}$. Derive the differential equation relating the displacement $x(t)$ to an external force $F(t)$.

\item Find the transfer function $G(s) = X(s)/F(s)$ for the system in part (a).

\item For initial conditions $x(0) = 0.01\text{ m}$ and $\dot{x}(0) = 0$, determine the system's response to a unit step input. Plot the response and identify the settling time and overshoot.

\end{enumerate}

% Section 1.3 Problems
\item \textbf{Thermal System Modeling}
\begin{enumerate}
\item Consider a heated room with thermal resistance $R_{th} = 0.5\text{ °C/W}$ and thermal capacitance $C_{th} = 5000\text{ J/°C}$. The heater provides a maximum power of $P = 2000\text{ W}$. Derive the differential equation for the room temperature $T(t)$ when the ambient temperature is $T_a = 10\text{ °C}$.

\item Find the transfer function relating room temperature to heater power input.

\item If the heater is turned on at full power at $t = 0$ with initial temperature $T(0) = 15\text{ °C}$, determine the steady-state temperature and the time constant of the system.

\end{enumerate}

% Section 1.4 Problems
\item \textbf{Translation to Rotation}
\begin{enumerate}
\item A DC motor is coupled to a load through a gear train with gear ratio $N = 10:1$ (motor turns 10 times for each load turn). The motor develops torque $T_m$ and the load requires torque $T_L$. Derive the equivalent inertia seen at the motor shaft, $J_{eq}$, in terms of load inertia $J_L$.

\item If the motor has armature resistance $R_a = 2\text{ \Omega}$, inductance $L_a = 0.01\text{ H}$, and torque constant $K_t = 0.1\text{ N·m/A}$, derive the transfer function relating angular velocity $\omega_m(s)$ to applied voltage $V_a(s)$.

\item Repeat part (b) for the case where back EMF constant $K_e = 0.1\text{ V·s/rad}$ is different from $K_t$. Discuss the implications.

\end{enumerate}

% Section 1.5 Problems
\item \textbf{State-Space Representation}
\begin{enumerate}
\item Convert the following differential equation to state-space form:
\[
\ddot{y} + 5\dot{y} + 6y = 3u
\]
where $u(t)$ is the input and $y(t)$ is the output. Choose the phase variable canonical form.

\item For the system in part (a), determine the matrices $A$, $B$, $C$, and $D$.

\item Find the transfer function from the state-space representation in part (b) and verify it matches the direct transfer function derivation.

\end{enumerate}

\item \textbf{Multivariable System}
\begin{enumerate}
\item A two-mass spring system has masses $m_1$ and $m_2$ connected by a spring $k$ to each other and to walls by springs $k_1$ and $k_2$. Define the state vector and derive the state-space equations with external forces $F_1$ and $F_2$ as inputs.

\item Write the state matrices $A$, $B$, $C$, and $D$ for the case $m_1 = m_2 = 1\text{ kg}$ and $k_1 = k_2 = k = 100\text{ N/m}$.

\end{enumerate}

% Section 1.6 Problems
\item \textbf{Block Diagram Manipulation}
\begin{enumerate}
\item Simplify the block diagram shown in Figure 1.XX and find the overall transfer function $Y(s)/R(s)$.

\item For a system with forward path transfer function $G(s) = 10/(s+2)(s+5)$ and feedback path transfer function $H(s) = 0.5s$, determine the closed-loop transfer function.

\item Determine the poles and zeros of the closed-loop system in part (b). Is the system stable?

\end{enumerate}

\item \textbf{Signal Flow Graph}
\begin{enumerate}
\item Convert the block diagram from Problem 1.15(a) to a signal flow graph using Mason's gain formula.

\item Apply Mason's gain formula to find the overall transfer function $Y(s)/R(s)$ for the signal flow graph.

\end{enumerate}

% Section 1.7 Problems
\item \textbf{Linearization}
\begin{enumerate}
\item Consider the nonlinear system $\dot{x} = x^2 + u$. Linearize this system about the equilibrium point $x_{eq} = 1$ with $u_{eq} = -1$.

\item For the linearized system in part (a), find the transfer function relating deviations $\tilde{x}(t)$ to deviations $\tilde{u}(t)$.

\item The pendulum equation is $\ddot{\theta} + \frac{g}{L}\sin\theta = \frac{1}{mL^2}u$. Linearize this equation about $\theta = \pi$ (the inverted equilibrium) and derive the linearized differential equation.

\end{enumerate}

\item \textbf{Numerical Linearization}
\begin{enumerate}
\item Consider the system $\dot{x} = \sin(x) + u$ with equilibrium at $x_{eq} = 0$. Derive the analytical linearization and compute the linearized transfer function.

\item Use numerical differentiation with $\Delta x = 0.001$ to approximate the linearization at $x_{eq} = 0$. Compare with the analytical result.

\end{enumerate}

% Section 1.8 Problems
\item \textbf{System Properties}
\begin{enumerate}
\item Determine whether each of the following systems is linear, time-invariant, causal, and stable:
\begin{enumerate}
\item $y(t) = t \cdot x(t)$
\item $y(t) = x(t-2)$
\item $y(t) = x^2(t)$
\item $y(t) = \int_{0}^{t} x(\tau)d\tau$
\end{enumerate}

\item For each system in part (a), provide a proof or counterexample supporting your conclusions.

\end{enumerate}

\item \textbf{Causality and Stability Analysis}
\begin{enumerate}
\item A system has impulse response $h(t) = e^{-2t}u(t) + e^{3t}u(-t)$. Is this system causal? Is it stable?

\item The transfer function of a system is $G(s) = \frac{s+1}{s^2 + 3s + 2}$. Determine if the system is BIBO stable by examining the pole locations.

\end{enumerate}

% Mixed Review Problems
\item \textbf{Comprehensive Modeling}
\begin{enumerate}
\item A liquid level system consists of a tank with cross-sectional area $A = 1\text{ m}^2$ and an outlet pipe with resistance $R = 2\text{ s/m}^2$. Derive the differential equation relating the liquid level $h(t)$ to the inflow rate $q_{in}(t)$.

\item Find the transfer function $H(s)/Q_{in}(s)$ for the system in part (a).

\item If a pump with flow rate $q_{in}(t) = 0.1\text{ m}^3/\text{s}$ is suddenly applied at $t = 0$, determine the steady-state liquid level and the time required to reach 63.2\% of the steady-state value.

\end{enumerate}

\item \textbf{Design and Selection}
\begin{enumerate}
\item \textbf{Design Problem:} You are tasked with modeling a pick-and-place robot for a control systems course project. The robot arm has three joints, each actuated by a DC motor.
\begin{enumerate}
\item What modeling approach would you recommend for the overall system: lumped parameter or distributed parameter? Justify your choice.
\item Select an appropriate level of detail for the motor models (include armature dynamics, neglect inductance, or use simplified model). Explain your reasoning.
\item Propose a state vector for the complete system that captures essential dynamics while remaining tractable for control design.
\end{enumerate}

\item \textbf{Design Problem:} A wind turbine manufacturer needs a model of the turbine-generator system for controller development. The system includes:
\begin{itemize}
\item Aerodynamic conversion from wind to mechanical torque
\item Shaft dynamics with possible flexibility
\item Generator dynamics (electrical and mechanical)
\item Pitch control for blade angle adjustment
\end{itemize}
\begin{enumerate}
\item What subsystems require detailed modeling versus simplified treatment?
\item Suggest transfer function forms for each major subsystem.
\item Identify potential coupling effects that might require multivariable modeling.

\end{enumerate}

\end{enumerate}

\item \textbf{Comparison and Analysis}
\begin{enumerate}
\item Compare and contrast the advantages and disadvantages of transfer function representations versus state-space representations for control system design. Provide examples where each representation is more suitable.

\item Discuss the conditions under which linearization of a nonlinear system produces a valid approximation of the original system's behavior. What factors affect the validity of the linear approximation?

\end{enumerate}

\item \textbf{Advanced Problem: Coupled Systems}
\begin{enumerate}
\item An electromechanical system consists of an electrical circuit (RLC network) mechanically coupled to a mass-spring-damper through a piezoelectric transducer. The transducer generates force proportional to voltage and voltage proportional to displacement.
\begin{enumerate}
\item Develop the coupled differential equations for this system.
\item Derive the transfer function from circuit input voltage to mass displacement.
\item Identify all coupling terms and discuss how they affect the system dynamics.

\end{enumerate}

\end{enumerate}

\item \textbf{Parameter Sensitivity}
\begin{enumerate}
\item For the mass-spring-damper system with $m = 1\text{ kg}$, $k = 100\text{ N/m}$, and $c = 2\text{ N·s/m}$:
\begin{enumerate}
\item Calculate the natural frequency, damping ratio, and damped natural frequency.
\item If the mass increases by 10\%, how does the natural frequency change? What is the percentage change?
\item Repeat for a 10\% increase in spring constant.

\end{enumerate}

\item Discuss how parameter uncertainty affects control system design and how sensitivity analysis can inform controller robustness requirements.

\end{enumerate}

\end{enumerate}

\bigskip

\textbf{Hints:}
\begin{itemize}
\item For Problems 1.1-1.3, remember to identify all energy storage elements as they determine the order of the dynamic model.
\item For Problems 1.9-1.10, the choice of state variables is not unique; verify your representation is controllable and observable when possible.
\item For linearization problems, always verify your Taylor series expansions.
 linearization by computing\item When simplifying block diagrams, carefully track how feedback signals combine at summing junctions.
\end{itemize}
